\section{Seguranca}

\subsection{Controle de Acesso}

\subsubsection{Roles e Permissoes}

\begin{lstlisting}[caption=Criacao de roles de seguranca]
-- Role para aplicacao (conexao principal)
CREATE ROLE finops_app WITH LOGIN PASSWORD 'senha_forte_aqui';

-- Role para backups
CREATE ROLE finops_backup WITH LOGIN PASSWORD 'senha_backup_aqui';

-- Role para leitura (relatorios)
CREATE ROLE finops_readonly WITH LOGIN PASSWORD 'senha_readonly_aqui';

-- Role para administracao
CREATE ROLE finops_admin WITH LOGIN SUPERUSER PASSWORD 'senha_admin_aqui';

-- Conceder permissoes basicas para aplicacao
GRANT CONNECT ON DATABASE finops TO finops_app;
GRANT USAGE ON SCHEMA public TO finops_app;
GRANT SELECT, INSERT, UPDATE, DELETE ON ALL TABLES IN SCHEMA public TO finops_app;
GRANT USAGE, SELECT ON ALL SEQUENCES IN SCHEMA public TO finops_app;
GRANT EXECUTE ON ALL FUNCTIONS IN SCHEMA public TO finops_app;

-- Permissoes para auditoria
GRANT USAGE ON SCHEMA audit TO finops_app;
GRANT INSERT ON audit.audit_logs TO finops_app;

-- Permissoes para readonly
GRANT CONNECT ON DATABASE finops TO finops_readonly;
GRANT USAGE ON SCHEMA public TO finops_readonly;
GRANT SELECT ON ALL TABLES IN SCHEMA public TO finops_readonly;

-- Permissoes para backup
GRANT CONNECT ON DATABASE finops TO finops_backup;
GRANT USAGE ON SCHEMA public, audit TO finops_backup;
GRANT SELECT ON ALL TABLES IN SCHEMA public, audit TO finops_backup;
\end{lstlisting}

\subsubsection{Row Level Security (RLS)}

\begin{lstlisting}[caption=Politicas de seguranca por linha]
-- Habilitar RLS nas tabelas principais
ALTER TABLE usuarios ENABLE ROW LEVEL SECURITY;
ALTER TABLE contas ENABLE ROW LEVEL SECURITY;
ALTER TABLE transacoes ENABLE ROW LEVEL SECURITY;
ALTER TABLE categorias ENABLE ROW LEVEL SECURITY;

-- Politica para usuarios (apenas dados proprios)
CREATE POLICY usuarios_self_access ON usuarios
    FOR ALL TO finops_app
    USING (id = current_setting('app.current_user_id')::uuid);

-- Politica para contas (proprias ou de grupo)
CREATE POLICY contas_access ON contas
    FOR ALL TO finops_app
    USING (
        usuario_id = current_setting('app.current_user_id')::uuid OR
        (compartilhada = true AND grupo_id IN (
            SELECT grupo_id FROM usuario_grupos 
            WHERE usuario_id = current_setting('app.current_user_id')::uuid 
              AND ativo = true
        ))
    );

-- Politica para transacoes
CREATE POLICY transacoes_access ON transacoes
    FOR ALL TO finops_app
    USING (
        usuario_id = current_setting('app.current_user_id')::uuid OR
        grupo_id IN (
            SELECT grupo_id FROM usuario_grupos 
            WHERE usuario_id = current_setting('app.current_user_id')::uuid 
              AND ativo = true
        )
    );

-- Politica para categorias
CREATE POLICY categorias_access ON categorias
    FOR ALL TO finops_app
    USING (
        sistema = true OR
        usuario_id = current_setting('app.current_user_id')::uuid OR
        grupo_id IN (
            SELECT grupo_id FROM usuario_grupos 
            WHERE usuario_id = current_setting('app.current_user_id')::uuid 
              AND ativo = true
        )
    );
\end{lstlisting}

\subsection{Criptografia}

\subsubsection{Criptografia de Dados Sensiveis}

\begin{lstlisting}[caption=Funcoes para criptografia]
-- Extensao para criptografia
CREATE EXTENSION IF NOT EXISTS pgcrypto;

-- Funcao para hash de senhas
CREATE OR REPLACE FUNCTION hash_senha(senha TEXT)
RETURNS TEXT AS $$
BEGIN
    RETURN crypt(senha, gen_salt('bf', 12));
END;
$$ LANGUAGE plpgsql;

-- Funcao para verificar senha
CREATE OR REPLACE FUNCTION verificar_senha(senha TEXT, hash TEXT)
RETURNS BOOLEAN AS $$
BEGIN
    RETURN hash = crypt(senha, hash);
END;
$$ LANGUAGE plpgsql;

-- Funcao para criptografar dados sensiveis
CREATE OR REPLACE FUNCTION criptografar_dado(dado TEXT, chave TEXT)
RETURNS TEXT AS $$
BEGIN
    RETURN encode(pgp_sym_encrypt(dado, chave), 'base64');
END;
$$ LANGUAGE plpgsql;

-- Funcao para descriptografar dados
CREATE OR REPLACE FUNCTION descriptografar_dado(dado_criptografado TEXT, chave TEXT)
RETURNS TEXT AS $$
BEGIN
    RETURN pgp_sym_decrypt(decode(dado_criptografado, 'base64'), chave);
EXCEPTION WHEN OTHERS THEN
    RETURN NULL;
END;
$$ LANGUAGE plpgsql;
\end{lstlisting}

\subsection{Auditoria e Logging}

\subsubsection{Log de Acessos}

\begin{lstlisting}[caption=Trigger para log de acessos]
CREATE TABLE audit.login_logs (
    id BIGSERIAL PRIMARY KEY,
    usuario_id UUID REFERENCES usuarios(id),
    ip_origem INET NOT NULL,
    user_agent TEXT,
    sucesso BOOLEAN NOT NULL,
    motivo_falha TEXT,
    timestamp_login TIMESTAMP DEFAULT CURRENT_TIMESTAMP
);

-- Funcao para log de login
CREATE OR REPLACE FUNCTION registrar_login(
    p_usuario_id UUID,
    p_ip INET,
    p_user_agent TEXT,
    p_sucesso BOOLEAN,
    p_motivo_falha TEXT DEFAULT NULL
)
RETURNS void AS $$
BEGIN
    INSERT INTO audit.login_logs (usuario_id, ip_origem, user_agent, sucesso, motivo_falha)
    VALUES (p_usuario_id, p_ip, p_user_agent, p_sucesso, p_motivo_falha);
END;
$$ LANGUAGE plpgsql;
\end{lstlisting}

\subsubsection{Deteccao de Anomalias}

\begin{lstlisting}[caption=Funcao para detectar atividades suspeitas]
CREATE OR REPLACE FUNCTION detectar_atividade_suspeita(p_usuario_id UUID)
RETURNS TABLE (
    tipo_alerta TEXT,
    descricao TEXT,
    severidade INTEGER
) AS $$
BEGIN
    -- Multiplos logins de IPs diferentes
    IF EXISTS (
        SELECT 1 FROM audit.login_logs 
        WHERE usuario_id = p_usuario_id 
          AND timestamp_login > CURRENT_TIMESTAMP - INTERVAL '1 hour'
          AND sucesso = true
        GROUP BY ip_origem 
        HAVING COUNT(*) > 5
    ) THEN
        RETURN QUERY SELECT 'MULTIPLOS_IPS'::TEXT, 'Multiplos logins de IPs diferentes'::TEXT, 3;
    END IF;
    
    -- Transacoes em valores muito altos
    IF EXISTS (
        SELECT 1 FROM transacoes t
        WHERE t.usuario_id = p_usuario_id
          AND t.data_transacao > CURRENT_DATE - INTERVAL '1 day'
          AND t.valor > (
              SELECT COALESCE(AVG(valor) * 10, 1000) 
              FROM transacoes t2 
              WHERE t2.usuario_id = p_usuario_id 
                AND t2.data_transacao > CURRENT_DATE - INTERVAL '30 days'
          )
    ) THEN
        RETURN QUERY SELECT 'VALOR_ALTO'::TEXT, 'Transacao com valor muito acima da media'::TEXT, 2;
    END IF;
    
    -- Muitas transacoes em pouco tempo
    IF (
        SELECT COUNT(*) FROM transacoes 
        WHERE usuario_id = p_usuario_id 
          AND created_at > CURRENT_TIMESTAMP - INTERVAL '10 minutes'
    ) > 20 THEN
        RETURN QUERY SELECT 'MUITAS_TRANSACOES'::TEXT, 'Muitas transacoes em pouco tempo'::TEXT, 2;
    END IF;
END;
$$ LANGUAGE plpgsql;
\end{lstlisting}

\subsection{Configuracoes de Seguranca}

\subsubsection{Parametros do PostgreSQL}

\begin{lstlisting}[caption=Configuracoes recomendadas de seguranca]
-- No postgresql.conf
-- ssl = on
-- ssl_cert_file = 'server.crt'
-- ssl_key_file = 'server.key'
-- ssl_ciphers = 'HIGH:MEDIUM:+3DES:!aNULL'
-- password_encryption = scram-sha-256
-- log_connections = on
-- log_disconnections = on
-- log_statement = 'mod'
-- log_line_prefix = '%t [%p]: [%l-1] user=%u,db=%d,app=%a,client=%h '

-- No pg_hba.conf
-- local   finops   finops_app               scram-sha-256
-- host    finops   finops_app  127.0.0.1/32 scram-sha-256
-- hostssl finops   finops_app  0.0.0.0/0    scram-sha-256
\end{lstlisting}

\subsection{Backup Seguro}

\begin{lstlisting}[caption=Script de backup com criptografia]
#!/bin/bash
# backup_seguro.sh

# Configuracoes
DB_NAME="finops"
BACKUP_DIR="/var/backups/finops"
DATE=$(date +%Y%m%d_%H%M%S)
BACKUP_FILE="$BACKUP_DIR/finops_backup_$DATE.sql.gz.gpg"

# Criar diretorio se nao existir
mkdir -p $BACKUP_DIR

# Realizar backup com compressao e criptografia
pg_dump $DB_NAME | gzip | gpg --cipher-algo AES256 --compress-algo 1 \
    --symmetric --output $BACKUP_FILE

# Verificar se backup foi criado
if [ -f "$BACKUP_FILE" ]; then
    echo "Backup criado com sucesso: $BACKUP_FILE"
    
    # Remover backups antigos (manter apenas 30 dias)
    find $BACKUP_DIR -name "finops_backup_*.sql.gz.gpg" -mtime +30 -delete
else
    echo "ERRO: Falha ao criar backup"
    exit 1
fi
\end{lstlisting}
